% Standalone: compile this file directly for the detailed theory supplement.
% Embedded: define \CGAChapterFiveEmbedded before inputting this file.
% In embedded mode the host provides theorem environments and bibliography.
\ifdefined\CGAChapterFiveEmbedded
\else
\documentclass[11pt]{article}
\usepackage[a4paper,margin=25mm]{geometry}
\usepackage{amsmath,amssymb,amsthm,mathtools,booktabs,microtype,enumitem}
\usepackage[hidelinks]{hyperref}
\usepackage{cleveref}
\def\CGAChapterFiveStandalone{1}
\newtheorem{theorem}{Theorem}[section]
\newtheorem{lemma}[theorem]{Lemma}
\newtheorem{proposition}[theorem]{Proposition}
\newtheorem{corollary}[theorem]{Corollary}
\newtheorem{assumption}[theorem]{Assumption}
\theoremstyle{remark}
\newtheorem{remark}[theorem]{Remark}
\setlength{\parskip}{3pt}
\emergencystretch=2em
\begin{document}
\setcounter{section}{3}
\fi
\providecommand{\cE}{\mathcal E}
\providecommand{\cD}{\mathcal D}
\providecommand{\cP}{\mathcal P}
\providecommand{\cK}{\mathcal K}
\providecommand{\cnorm}[1]{\left\lVert#1\right\rVert}
\providecommand{\cabs}[1]{\left|#1\right|}
\providecommand{\cinner}[2]{\left\langle#1,#2\right\rangle}
\providecommand{\cdd}{\,\mathrm d}
\providecommand{\cR}{\mathbb R}

\ifdefined\CGAChapterFiveStandalone
  \newcommand{\CGAHLDreference}{the finite-pool descent estimates}
  \newcommand{\CGAEnergyMetricReference}{the energy--distance theorem and metric corollary in the main text}
  \newcommand{\CGAEnergyTheoremReference}{the energy--distance theorem in the main text}
\else
  \newcommand{\CGAHLDreference}{Section~\ref{sec:descent}}
  \newcommand{\CGAEnergyMetricReference}{Theorem~\ref{thm:energy-v} and Corollary~\ref{cor:metric-rates}}
  \newcommand{\CGAEnergyTheoremReference}{Theorem~\ref{thm:energy-v}}
\fi

\section{Convergence analysis and finite-trajectory projection bounds}
\label{sec:c5}

Finite-pool descent accounts for continuous-dictionary coverage, score evaluation, radial stationarity, and coefficient correction. Local Hessian transfer connects nonlinear selection with weak OGA in a frozen Hilbert metric. Projection estimates describe the nested spaces generated by the algorithm, and the quartic identity converts their approximation errors into energy bounds.

Throughout, \(m\) counts accepted independent directions, \(m_0\) is the starting index, and \(j=m-m_0\). We use \(\widehat G_m\) for the computed state, \(G_m^\sharp\) for the exact active-space energy minimizer, and \(z_m\) for the frozen orthogonal projection.
All derivatives below act on the specified feasible linear space, so that
$\cE'(u^*)=0$ there. Affine constraints are reduced to this setting by a
fixed feasible lift.

\subsection{Finite-pool descent and the residual-comparison condition}
\label{sec:c5-descent}

Let \(X\) be a real Banach space and \(\cE:X\to(-\infty,+\infty]\) a proper convex energy with minimizer \(u^*\). Assume that \(\cE\) is finite and continuously differentiable on an open neighborhood containing the states and dictionary segments used below.
The symmetric dictionary $\cD_X$ consists of $X$-unit vectors. Set
\[
 \widehat\delta_m=\cE(\widehat G_m)-\cE(u^*),\qquad
 S_m^X=\sup_{g\in\cD_X}\cabs{\cinner{\cE'(\widehat G_m)}g}.
\]
The exact finite-pool score is $s_m(g)=|\langle\cE'(\widehat G_m),g\rangle|$.
The available computed scores satisfy
\begin{equation}
 \cabs{\widetilde s_m(g)-s_m(g)}\leq\xi_m\quad(g\in\cP_m),\qquad
 \widetilde s_m(g_{m+1})\geq t_m^{\rm pool}
                  \max_{g\in\cP_m}\widetilde s_m(g),
 \label{eq:c5-oracle}
\end{equation}
where $\cP_m\subset\cD_X$ is finite and $0<t_m^{\rm pool}\leq1$.
When $S_m^X>0$, define its continuous coverage by
\begin{equation}
 \tau_m=\frac{\max_{g\in\cP_m}s_m(g)}{S_m^X}\in[0,1].
 \label{eq:c5-coverage}
\end{equation}
Set $\tau_m=1$ when $S_m^X=0$. The selection weakness $t_m^{\rm pool}$
and the coverage $\tau_m$ are different quantities.

Assume that a common $L>0$ satisfies
\begin{equation}
 \cE(\widehat G_m+\lambda g)
 \leq\cE(\widehat G_m)+\lambda\langle\cE'(\widehat G_m),g\rangle
            +\tfrac L2\lambda^2,\qquad |\lambda|\leq1,\quad g\in\cD_X.
 \label{eq:c5-directional}
\end{equation}
The expanded space contains $\widehat G_m+\lambda g_{m+1}$, and the
correction satisfies
\begin{equation}
 \cE(\widehat G_{m+1})\leq
 \inf_{v\in\widehat\Phi_{m+1}}\cE(v)+\varepsilon_m^{\rm opt},
 \qquad \varepsilon_m^{\rm opt}\geq0 .
 \label{eq:c5-correction}
\end{equation}
The active-space energy comparison relates this correction estimate to the numerical objective.

\begin{proposition}[Descent with a resolved score]
\label{prop:c5-descent}
Let $a_m\geq0$ be any valid lower bound for $s_m(g_{m+1})$.
Under \eqref{eq:c5-directional}--\eqref{eq:c5-correction},
\begin{equation}
 \widehat\delta_{m+1}\leq\widehat\delta_m-H_L(a_m)
                       +\varepsilon_m^{\rm opt}.
 \label{eq:c5-descent}
\end{equation}
Here $H_L$ is the directional-decrease function defined in \CGAHLDreference.
In particular, \eqref{eq:c5-oracle} gives the valid lower bound
\begin{equation}
 a_m=\left[t_m^{\rm pool}\tau_m S_m^X
                 -(1+t_m^{\rm pool})\xi_m\right]_+ .
 \label{eq:c5-score-lower}
\end{equation}
\end{proposition}
\begin{proof}
The triangle inequality gives
$s_m(g_{m+1})\geq
t_m^{\rm pool}\max_{\cP_m}s_m-(1+t_m^{\rm pool})\xi_m$.
Test the correction with a signed dictionary step in the true descent
direction of length $\min\{1,a_m/L\}$. Both signs lie in the expanded
linear space. Minimizing the quadratic upper bound on that segment
gives $H_L(a_m)$ and proves \eqref{eq:c5-descent}.
\end{proof}

The radial stationarity error enters when an atomic source is used to
bound $S_m^X$. Write $K_1(\cD)$ for the closed absolutely convex hull
of a bounded dictionary. If
$u^*\in B_{\rm src}K_1(\cD_X)$ with $B_{\rm src}>0$ and
$|\langle\cE'(\widehat G_m),\widehat G_m\rangle|
\leq\eta_m^{\rm stat}$, convexity implies
\begin{equation}
 \widehat\delta_m\leq\eta_m^{\rm stat}+B_{\rm src}S_m^X .
 \label{eq:c5-atomic}
\end{equation}
Consequently one may instead use
\[
 a_m=\left[
 t_m^{\rm pool}\tau_m
       \frac{\widehat\delta_m-\eta_m^{\rm stat}}{B_{\rm src}}
 -(1+t_m^{\rm pool})\xi_m\right]_+ .
\]
This source-dependent estimate expresses the selected-score lower bound in terms of the energy gap and radial stationarity.
In the error framework of \cite[Section~3, Theorem~3.2 and
Lemma~6.3]{DereventsovTemlyakov2022}, energy correction and radial
biorthogonality are separate errors. Their power-smoothness rate assumes
uniform positive continuous weakness and prescribed decay of those errors.
Here they correspond to $\varepsilon_m^{\rm opt}$ and
$\eta_m^{\rm stat}$, with the latter at the preceding state in
\eqref{eq:c5-atomic}. Pool coverage and score evaluation determine the effective selection strength through \eqref{eq:c5-score-lower}.

\begin{lemma}[A nonautonomous scalar bound]
\label{lem:c5-recurrence}
Suppose $x_m\geq0$ and
$x_{m+1}\leq x_m-b_mx_m^q$ for $q>1$ and $b_m\geq0$.
For $x_0>0$,
\begin{equation}
 x_n\leq\left[x_0^{1-q}+(q-1)\sum_{m=0}^{n-1}b_m\right]^{-1/(q-1)} .
 \label{eq:c5-recurrence}
\end{equation}
If $x_0=0$, the sequence vanishes. In particular,
$b_m\geq b(m+1)^{-s}$, $b>0$, $0\leq s<1$, implies
$x_n=O((n+1)^{-(1-s)/(q-1)})$.
\end{lemma}
\begin{proof}
As long as $x_{m+1}>0$, convexity gives
$x_{m+1}^{1-q}-x_m^{1-q}\geq
(q-1)x_m^{-q}(x_m-x_{m+1})\geq(q-1)b_m$.
Sum. Once a term vanishes, nonnegativity and the recurrence force all
later terms to vanish. The power estimate follows by comparison of the
sum with an integral.
\end{proof}

\begin{theorem}[Residual-comparison rate]
\label{thm:c5-visible}
Let $p>2$. Suppose, at the computed states,
\begin{align}
 \widehat\delta_m&\geq c_U\cnorm{\widehat G_m-u^*}_X^p,\quad c_U>0,
 \label{eq:c5-uniform-convex}\\
 S_m^X&\geq\kappa_0(m+1)^{-\chi}
                   \cnorm{\cE'(\widehat G_m)}_{X^*},
 \quad\kappa_0>0,\quad0\leq\chi<\tfrac12 .
 \label{eq:c5-visible}
\end{align}
Assume \eqref{eq:c5-oracle}--\eqref{eq:c5-correction}, and, uniformly in $m$,
\[
 t_m^{\rm pool}\tau_m\geq\underline t>0,\qquad
 (1+t_m^{\rm pool})\xi_m
 \leq\vartheta t_m^{\rm pool}\tau_m S_m^X,\quad0\leq\vartheta<1 .
\]
With $a_m$ as in \eqref{eq:c5-score-lower}, assume also
$\varepsilon_m^{\rm opt}\leq\zeta H_L(a_m)$, $0\leq\zeta<1$.
Then
\begin{equation}
 \widehat\delta_m=O((m+1)^{-p(1-2\chi)/(p-2)}) .
 \label{eq:c5-visible-rate}
\end{equation}
The constant depends on $p,c_U,\kappa_0,\underline t,L,\vartheta,\zeta$ and
$\widehat\delta_0$. For the exact continuous method take
$\tau_m=1$, $\xi_m=\varepsilon_m^{\rm opt}=0$.
\end{theorem}
\begin{proof}
Convexity and \eqref{eq:c5-uniform-convex} give
\[
 \widehat\delta_m\leq
 \cnorm{\cE'(\widehat G_m)}_{X^*}
       (\widehat\delta_m/c_U)^{1/p}.
\]
For positive error this yields gradient domination.
Thus $a_m\geq b_m^*:=K(m+1)^{-\chi}
\widehat\delta_m^{(p-1)/p}$, where
$K=(1-\vartheta)\underline t\kappa_0c_U^{1/p}$.
Descent and absorption imply monotonicity, so
$b_m^*\leq A:=K\widehat\delta_0^{(p-1)/p}$.
For $0\leq b\leq A$,
$H_L(b)\geq b^2/(2\max\{L,A\})$.
Since $H_L$ is increasing, \eqref{eq:c5-descent} implies
\[
 \widehat\delta_{m+1}\leq\widehat\delta_m-
 \frac{(1-\zeta)K^2}{2\max\{L,A\}}
 (m+1)^{-2\chi}\widehat\delta_m^{2(p-1)/p}.
\]
Apply Lemma~\ref{lem:c5-recurrence}. The case $\widehat\delta_0=0$
is immediate from descent and absorption.
\end{proof}

The norm comparison is imposed along the specified residual trajectory. The resulting energy exponent exceeds the atomic baseline exponent \(1\) when \(\chi<1/p\).

\subsection{Energy and the natural distance}
\label{sec:c5-natural}

For the standard pure, regularized, and reaction $p$-energies, the energy--distance
equivalence and the $W^{1,p}$ conversion are established in
\CGAEnergyMetricReference. Therefore, if the
energy estimate in Theorem~\ref{thm:c5-visible} is
$\widehat\delta_m=O(m^{-\alpha})$, then
\[
 d_{V,\bullet}(\widehat G_m,u^*)=O(m^{-\alpha/2}),\qquad
 \|\widehat G_m-u^*\|_{W^{1,p}}=O(m^{-\alpha/p}),
\]
under the corresponding zero-mean or reaction assumptions.

Thus \eqref{eq:c5-visible-rate} has natural-distance exponent \(p(1-2\chi)/(2(p-2))\) and Sobolev exponent \((1-2\chi)/(p-2)\) for the pure, regularized, and reaction energies.

\subsection{Local Hessian transfer}
\label{sec:c5-hessian}

The transfer argument connects nonlinear greedy selection with Hilbert-space approximation. Xu--Xu \cite[Theorem~3.1, equations~(17)--(19)]{XuXu2025} analyze exact active-space minimization of a convex, $L$-smooth functional in a Sobolev Hilbert space. In their notation, the continuous
selection strengths $\gamma_i$, source radius $B$, and innovations
$\bar g_i=(I-P_{i-1})g_i$ yield
\[
 \delta_i\leq\delta_{i-1}
 -\frac{\gamma_i^2\delta_{i-1}^2}{2LB^2\cnorm{\bar g_i}^2}.
\]
Their ReLU dictionary estimate then gives an energy rate
$n^{-1-(2(k-s)+1)/d}$ for constant positive weakness in $H^s$,
with $k\geq s$; strong convexity gives the corresponding norm bound.
Here \(s\) denotes the Sobolev order. The argument below transfers nonlinear selection to a frozen metric through explicit local Hessian and relative-score conditions. The finite-pool estimate separates coverage, score evaluation, radial stationarity, and coefficient correction. The entropy argument applies the weak-selection volume estimate to the projected atoms with their inherited norms.

Let $H$ be a real Hilbert space on which $\cE$ is twice Fr\'echet
differentiable in an open convex neighborhood of $u^*$.
Assume $A_*=\cE''(u^*)$ is bounded, symmetric, and coercive on $H$.
Write $a_*(v,w)=\langle A_*v,w\rangle$ and $\cnorm v_*^2=a_*(v,v)$.
Suppose the closed ball of radius $R_*$ around $u^*$ lies in that
neighborhood and
\begin{equation}
 |\langle[\cE''(u)-A_*]v,w\rangle|
 \leq L_*\cnorm{u-u^*}_*\cnorm v_*\cnorm w_*,
 \qquad \cnorm{u-u^*}_*\leq R_*,\quad L_*R_*\leq\tfrac12 .
 \label{eq:c5-hessian}
\end{equation}
For finite-dimensional $\Phi\subset H$ let $z_\Phi$ be the frozen
orthogonal projection of $u^*$ and $G_\Phi^\sharp$ an exact minimizer
of $\cE$ over $\Phi$. Assume both points lie in this ball.

\begin{lemma}[Taylor and projection estimates]
\label{lem:c5-projection}
Under \eqref{eq:c5-hessian}, for $u$ in the ball,
\begin{align}
 \cabs{\cE(u)-\cE(u^*)-\tfrac12\cnorm{u-u^*}_*^2}
 &\leq\tfrac{L_*}{6}\cnorm{u-u^*}_*^3,\label{eq:c5-taylor}\\
 \cnorm{\cE'(u)-A_*(u-u^*)}_{*'}
 &\leq\tfrac{L_*}{2}\cnorm{u-u^*}_*^2,\label{eq:c5-grad-taylor}\\
 \cnorm{G_\Phi^\sharp-z_\Phi}_*
 &\leq L_*\cnorm{z_\Phi-u^*}_*^2.\label{eq:c5-proj}
\end{align}
Here the dual norm in \eqref{eq:c5-grad-taylor} is taken over $H$.
\end{lemma}
\begin{proof}
One and two integrations of \eqref{eq:c5-hessian} give the first
two estimates. Put $d=G_\Phi^\sharp-z_\Phi\in\Phi$.
Use $\langle\cE'(G_\Phi^\sharp),d\rangle=0$ and
$a_*(z_\Phi-u^*,d)=0$. Subtracting $A_*d$ along the segment from
$z_\Phi$ to $G_\Phi^\sharp$ bounds
$\cnorm d_*^2$ by
$L_*R_*\cnorm d_*^2+(L_*/2)\cnorm{z_\Phi-u^*}_*^2\cnorm d_*$.
Absorb the first term. All segments remain in the ball by convexity.
\end{proof}

Assume $\cD_X\subset H$ and the algorithmic dictionary satisfies the uniform norm comparison
\begin{equation}
 0<c_*\leq\cnorm g_*\leq C_*<\infty\quad(g\in\cD_X),\qquad
 \cD_*=\{g/\cnorm g_*:g\in\cD_X\}.
 \label{eq:c5-bridge}
\end{equation}
Let $\Phi_m$ be nested finite-dimensional spaces with
$\Phi_{m+1}=\Phi_m+\operatorname{span}\{g_{m+1}\}$, where each selected
direction is independent of $\Phi_m$. Denote their frozen orthogonal
projections by $P_m$, and set $z_m=P_mu^*$,
$r_m=\cnorm{u^*-z_m}_*$, and
$G_m^\sharp=G_{\Phi_m}^\sharp$.
At an exact active-space state set
\[
 S_m^*=\sup_{d\in\cD_*}|a_*(z_m-u^*,d)|,\qquad
 \eta_m^H=\sup_{g\in\cD_X}
 |\langle\cE'(G_m^\sharp),g\rangle-a_*(z_m-u^*,g)|.
\]
The preceding lemma gives \(\eta_m^H\le C_T\cnorm{z_m-u^*}_*^2\), uniformly on the specified ball and dictionary. The relative-score condition below compares this defect with \(S_m^*\).

\begin{proposition}[Selection under a relative score condition]
\label{prop:c5-transfer}
Suppose nonlinear selection over $\cD_X$ has strength $t\in(0,1]$.
For $0<t_0<tc_*/C_*$, the condition
\begin{equation}
 (1+t)\eta_m^H\leq(tc_*-t_0C_*)S_m^*
 \label{eq:c5-relative}
\end{equation}
implies
$|a_*(z_m-u^*,g_{m+1}/\cnorm{g_{m+1}}_*)|\geq t_0S_m^*$.
\end{proposition}
\begin{proof}
The selected unnormalized frozen score is at least
$t\sup_{g\in\cD_X}|a_*(z_m-u^*,g)|-(1+t)\eta_m^H$.
The supremum is at least $c_*S_m^*$, and division by the selected
norm costs at most $C_*$. This proves the assertion.
\end{proof}

Fix $m_0$, $Q_0=I-P_{m_0}$ and
\[
 H_0=Q_0H,\qquad f_0=Q_0u^*,\qquad
 \cD_{\rm rst}=Q_0\cD_*\setminus\{0\}.
\]
The projected atoms retain their projected norms, which are at most one. Under the \(2^{n-1}\)-ball convention, the \(2^j\)-ball cover in the volume argument corresponds to \(\varepsilon_{j+1}\).

\begin{lemma}[Entropy bound for bounded-atom weak OGA]
\label{lem:c5-entropy}
Let $\cD_0$ be symmetric with atom norms at most one in a Hilbert space,
$K=K_1(\cD_0)$ compact, and $f_0\in B_*K$ with $B_*\geq0$.
Let weak OGA of strength $t_0>0$ start from the zero space.
After $j\geq1$ independent extensions its error satisfies
\begin{equation}
 \cnorm{f_0-P_jf_0}\leq
 \frac{B_*}{t_0}\frac{(j!V_j)^{1/j}}{\sqrt j}
                    \varepsilon_{j+1}(K),\qquad
 V_j=\frac{\pi^{j/2}}{\Gamma(j/2+1)} .
 \label{eq:c5-entropy}
\end{equation}
\end{lemma}
\begin{proof}
If $B_*=0$, then $f_0=0$ and the assertion is immediate. Suppose $B_*>0$.
Write $y_i=\cnorm{f_0-P_if_0}^2$ and
$l_i=\cnorm{(I-P_{i-1})d_i}>0$.
Orthogonality and the source condition give the best score at least
$y_{i-1}/B_*$. Hence
$y_i\leq y_{i-1}-t_0^2y_{i-1}^2/(B_*^2l_i^2)$.
If the error reaches zero the result is immediate. Otherwise taking
reciprocals and using the arithmetic--geometric mean inequality yields
\[
 y_j\leq \frac{B_*^2}{t_0^2\sum_{i=1}^jl_i^{-2}}
 \leq \frac{B_*^2}{t_0^2j}\Big(\prod_{i=1}^jl_i\Big)^{2/j}.
\]
The cross-polytope generated by $\{\pm d_1,\ldots,\pm d_j\}$ has
$j$-dimensional volume $(2^j/j!)\prod_i l_i$ and lies in the orthogonal
projection of $K$ onto their span. A cover of $K$ by $2^j$ balls of radius
\(\varepsilon_{j+1}(K)+h\), \(h>0\), projects to a cover of that
cross-polytope. Volume comparison gives
$\prod_i l_i\leq j!V_j(\varepsilon_{j+1}(K)+h)^j$.
Let $h\downarrow0$ and substitute.
\end{proof}

The proof adapts the volume--entropy argument of Li--Siegel \cite[Theorem~4.2 and equations~(4.4)--(4.9)]{LiSiegelEntropy2024} to weak selection with atom norms bounded by one, retaining the innovation weights of the restarted spaces. Stirling's formula gives a uniform bound on the factorial--volume factor divided by $\sqrt j$.

\begin{theorem}[Hessian transfer after restart]
\label{thm:c5-hessian}
Assume \eqref{eq:c5-hessian} and \eqref{eq:c5-bridge}.
For every state in $m_0\leq m\leq m_0+J$, suppose that the exact nonlinear
state and frozen projection lie in the specified ball.
For each selection step $m_0\leq m<m_0+J$, assume nonlinear selection
has strength at least $t$, and \eqref{eq:c5-relative} holds with
the same $t_0>0$. On the unrenormalized restarted dictionary assume
\begin{equation}
 f_0\in B_*K_1(\cD_{\rm rst}),\qquad
 \varepsilon_j(K_1(\cD_{\rm rst});H_0)\leq C_{\rm ent}j^{-\beta},
 \label{eq:c5-source-entropy}
\end{equation}
for $1\leq j\leq J+1$, with compact closed absolutely convex hull,
$B_*,C_{\rm ent}>0$, and $\beta>0$.
Then
\[
 \cnorm{z_{m_0+j}-u^*}_*\leq
 C B_*C_{\rm ent}t_0^{-1}j^{-\beta},\qquad
 \cE(G_{m_0+j}^\sharp)-\cE(u^*)\leq C_{\rm nl}j^{-2\beta}.
\]
Here \(C\) is absolute and \(C_{\rm nl}\) depends on \(L_*,R_*,B_*,C_{\rm ent},t_0\). The \(2^{n-1}\)-ball convention gives the index \(j+1\) in \Cref{lem:c5-entropy}, so the entropy assumption is stated through \(J+1\).
\end{theorem}
\begin{proof}
Proposition~\ref{prop:c5-transfer} gives weak frozen selection.
The residual is orthogonal to the starting space, so its scores against
$d$ and $Q_0d$ agree. The spaces after the starting index therefore form weak OGA
on $H_0$ with dictionary $\cD_{\rm rst}$. Apply
Lemma~\ref{lem:c5-entropy}. Lemma~\ref{lem:c5-projection} gives
$\cnorm{G_{m_0+j}^\sharp-u^*}_*
\leq r_{m_0+j}+L_*r_{m_0+j}^2
\leq(1+L_*R_*)r_{m_0+j}$, and \eqref{eq:c5-taylor}
gives the energy bound.
\end{proof}

\paragraph{A family satisfying the structural assumptions.}
Take $H=H^1(0,1)$ and
\[
 \cE_\lambda(u)=\int_0^1
 \left(\tfrac12|u'|^2+\tfrac12u^2+\tfrac\lambda4u^4\right)\cdd x-\ell(u),
 \qquad\lambda\geq0 .
\]
For a fixed integer $k\geq2$ and $0<a<b<1$, put
$g_t=(x-t)_+^k/\cnorm{(x-t)_+^k}_{H^1}$, $t\in[a,b]$,
and $\cD_X=\{\pm g_t\}$.
Choose any finite signed measure $\mu$ with total variation at most
$B_{\rm src}$, set $u^*=\int_a^b g_t\cdd\mu(t)$, and manufacture
$\ell(v)=\int(u^{*\prime}v'+u^*v+\lambda(u^*)^3v)$.
The embedding \(H^1(0,1)\hookrightarrow L^\infty\) makes the load continuous and the Hessian locally Lipschitz on \(H^1\) balls. The variational load determines the natural boundary flux.
Its bilinear form is
$\int(v'w'+(1+3\lambda(u^*)^2)vw)$, which is coercive in $H^1$.
The norm comparison \eqref{eq:c5-bridge} follows, the dictionary is compact,
and $u^*\in C_*B_{\rm src}K_1(\cD_*)$.
Its closed absolutely convex hull is compact: finite nets for the
dictionary approximate that hull by bounded finite-dimensional convex hulls.
Projection by \(Q_0\) preserves the source representation and bounds the projected entropy by the original entropy. These structural assumptions hold for the family of targets defined above.
For \(\lambda=0\), \(\eta_m^H=0\) and the setting reduces to Hilbert OGA. For \(\lambda>0\), the relative-score condition is expressed through the following quantitative small-nonlinearity estimate. Let
$C_{\rm emb}=\|H^1(0,1)\hookrightarrow L^\infty(0,1)\|$ (for the
standard squared $H^1$ norm one may use the explicit bound
$C_{\rm emb}\leq\sqrt2$) and
$M_* = \|u^*\|_{L^\infty}$, and use the $H^1$-normalized atoms above.
Since $a_*$ dominates the standard $H^1$ inner product, the quartic
Hessian satisfies, on the specified $R_*$-ball,
\begin{equation}
 L_\lambda:=3\lambda C_{\rm emb}
       (C_{\rm emb}R_*+2M_*),\qquad L_\lambda R_*\leq\tfrac12,
 \label{eq:c5-family-hessian-constant}
\end{equation}
and the norm-comparison and source constants can be chosen as
\begin{equation}
 c_*=1,\qquad C_\lambda=\sqrt{1+3\lambda M_*^2},\qquad
 B_\lambda=C_\lambda B_{\rm src}.
 \label{eq:c5-family-bridge-constants}
\end{equation}
Indeed, $\|g_t\|_*=\sqrt{1+3\lambda\|u^*g_t\|_{L^2}^2}$ lies
between $1$ and $C_\lambda$, and the change of variables from $g_t$ to
$g_t/\|g_t\|_*$ multiplies the atomic budget by at most $C_\lambda$.
Put
\begin{equation}
 C_{T,\lambda}=L_\lambda\left[1+\tfrac12(1+L_\lambda R_*)^2\right].
 \label{eq:c5-family-transfer-constant}
\end{equation}
The projection estimates then give
$\eta_m^H\leq C_\lambda C_{T,\lambda}r_m^2$.  Consequently the
following explicit bound is sufficient for the relative-score condition:
\begin{equation}
 \Delta_\lambda:=(tc_*-t_0C_\lambda)
 -(1+t)C_\lambda^2C_{T,\lambda}B_{\rm src}>0.
 \label{eq:c5-small-lambda-score}
\end{equation}
For fixed \(R_*\), \(M_*\), \(B_{\rm src}\), \(t\), and \(t_0<tc_*\), continuity of the nonlinear term as \(\lambda\downarrow0\) gives an interval \(0<\lambda\le\lambda_0\) satisfying \eqref{eq:c5-family-hessian-constant} and \eqref{eq:c5-small-lambda-score}.

\paragraph{Sufficient conditions for the concrete semilinear family.}
Fix \(m_0\), \(Q_0=I-P_{m_0}\), and the restarted dictionary \(\cD_{\rm rst}=Q_0\cD_*\setminus\{0\}\), with the projected norms retained and \(\cD_*=\{g_t/\|g_t\|_*:t\in[a,b]\}\) including both signs. The sufficient conditions on \(0\le j\le J\) are:
\begin{enumerate}[label=(A\arabic*)]
  \item an initial radius and a specified $R_*$ for which every exact active
  minimizer, frozen projection, and corrected state in the stated finite range stays in
  $B_*(u^*,R_*)$; for a specified $r_{\rm ent}>0$ this includes
  $r_{m_0}\leq r_{\rm ent}$ and
  $L_\lambda R_*\leq\tfrac12$;
  \item a uniform selection strength $t$ and $t_0<tc_*/C_\lambda$ for
  which the displayed bound $\Delta_\lambda$ is positive at every
  selection step $m_0\leq m<m_0+J$ in the stated range;
  \item a quantitative entropy estimate for this same projected,
  unrenormalized object, with the ball-counting convention fixed above,
  \begin{equation}
    \varepsilon_n\!\left(K_1(\cD_{\rm rst});H_0\right)
       \leq C_{\rm ent}n^{-\beta_A},
       \qquad 1\leq n\leq J+1,
    \label{eq:c5-family-entropy-certificate}
  \end{equation}
  and $Q_0u^*\in B_\lambda K_1(\cD_{\rm rst})$; and
  \item correction, quadrature, and radial errors absorbed in the
  finite-pool score bound, with tolerances prescribed from the available finite-pool quantities.
\end{enumerate}
The exponent \(\beta_A\) and constant \(C_{\rm ent}\) are those of the projected-dictionary entropy estimate \eqref{eq:c5-family-entropy-certificate}.

\begin{corollary}[Concrete-family rate]
\label{cor:c5-concrete-family}
Assume the sufficient conditions above hold for some $\lambda\in(0,\lambda_0]$
and $0<\beta_A$, and use exact active-space minimization in the statement
of the rate.  Then the restarted CGA trajectory satisfies, for
$1\leq j\leq J$,
\begin{equation}
 \cnorm{z_{m_0+j}-u^*}_*
   \leq C B_\lambda C_{\rm ent}t_0^{-1}j^{-\beta_A},
 \qquad
 \cE_\lambda(G_{m_0+j}^\sharp)-\cE_\lambda(u^*)
   \leq C_{\rm nl}j^{-2\beta_A}.
 \label{eq:c5-concrete-family-rate}
\end{equation}
The constants are those of Theorem~\ref{thm:c5-hessian}.
\end{corollary}

\paragraph{Function-space setting.}
The concrete semilinear family uses the compact parameter interval \([a,b]\subset(0,1)\). The sinh model uses an \(L^\infty\)-bounded neighborhood for local Hessian estimates and finite-dimensional spaces for the trajectory analysis. Pure \(p\)-energies are treated through the global energy geometry and the projection bounds below.

\subsection{Finite-reference scores and restricted stationarity}
\label{sec:c5-validation}

On a specified finite-dimensional space $Y_M$ with a positive frozen inner
product, assume $u^*\in Y_M$ and $\widehat\Phi_m\subset Y_M$.
In this setting let $z_m$ be the frozen projection of $u^*$ onto
$\widehat\Phi_m$. Let the fixed reference set $\mathcal A_M\subset Y_M$
contain all pools on the stated finite range.
Assume $0<c_M\leq\kappa_g:=\cnorm g_*\leq C_M<\infty$.
Set
\[
 h_m(g)=|a_*(u^*-z_m,g)|,\quad
 S_m^{\rm val}=\max_{g\in\mathcal A_M}h_m(g)/\kappa_g,\quad
 \tau_m^{\rm val}
   =\frac{\max_{g\in\cP_m}h_m(g)/\kappa_g}{S_m^{\rm val}}
\]
when $S_m^{\rm val}>0$.
If the computed scores obey
$|\widetilde s_m(g)-h_m(g)|\leq\Delta_m$ on $\cP_m$, the proof of
Proposition~\ref{prop:c5-transfer} gives
\begin{equation}
 \frac{h_m(g_{m+1})}{\kappa_{g_{m+1}}}
 \geq
 \frac{t_m^{\rm pool}c_M\tau_m^{\rm val}S_m^{\rm val}
             -(1+t_m^{\rm pool})\Delta_m}{C_M}.
 \label{eq:c5-finite-margin}
\end{equation}
Thus the following nonnegative lower bound
$(t_m^{\rm pool}c_M\tau_m^{\rm val}-t_0C_M)S_m^{\rm val}
 -(1+t_m^{\rm pool})\Delta_m$ suffices for weakness $t_0$
relative to this finite reference set.
The additional comparison \(S_m^{\rm val}\ge(1-\zeta_M)S_m^*\), with \(0\le\zeta_M<1\), transfers the finite-reference estimate to the continuous score. The numerical records report quadrature sensitivity through the comparison of two quadrature rules.

The active-space KKT residual in a reference norm $\cnorm\cdot_{0,M}$ is
\begin{equation}
 k_m=\sup_{\substack{v\in\widehat\Phi_m\\\cnorm v_{0,M}=1}}
              |\langle\cE'(\widehat G_m),v\rangle|.
 \label{eq:c5-kkt}
\end{equation}
The supremum is over the active space, even when its norm is inherited
from $Y_M$; set $k_m=0$ if $\widehat\Phi_m=\{0\}$.
If $\cE$ is $\mu_{\rm act}>0$ strongly convex on the active
space and its exact minimizer exists, then
\[
 \cnorm{\widehat G_m-G_m^\sharp}_{0,M}\leq k_m/\mu_{\rm act},\quad
 \cE(\widehat G_m)-\cE(G_m^\sharp)\leq k_m^2/(2\mu_{\rm act}).
\]
The first follows by strong monotonicity paired with the state difference;
the second follows by minimizing the strong-convexity quadratic.
A bound $C_{\rm score}$ on
$\cnorm{\cE''(u)g|_{\widehat\Phi_m}}_{\widehat\Phi_m^*}$ over the
relevant segments and pool gives an additional score error at most
$C_{\rm score}k_m/\mu_{\rm act}$.
The coercivity constants depend on the chosen reference norm and trial space. Since $\widehat\Phi_m\subset Y_M$,
\[
 \inf_{0\ne v\in\widehat\Phi_m}\frac{a_*(v,v)}{\cnorm v_{0,M}^2}
 \ \geq\
 \inf_{0\ne v\in Y_M}\frac{a_*(v,v)}{\cnorm v_{0,M}^2}.
\]
Thus the coercivity constant is associated with the space over which the Rayleigh quotient is minimized.

\subsection{An approximate-source estimate for the computed trajectory}
\label{sec:c5-source}

Let $Y$ be a finite-dimensional real Hilbert space containing $u^*$,
all active spaces, and a fixed finite symmetric dictionary
$\cD_{\rm fin}$ of frozen-unit vectors. For $m_0\leq m<m_0+J$ assume
\begin{equation}
 \widehat\Phi_{m+1}=\widehat\Phi_m+\operatorname{span}\{d_{m+1}\},
 \quad d_{m+1}\in\cD_{\rm fin},\quad d_{m+1}\notin\widehat\Phi_m .
 \label{eq:c5-spaces}
\end{equation}
The dictionary must contain every accepted direction in this finite range.
Here $P_m$ denotes the frozen orthogonal projection onto $\widehat\Phi_m$.
Set
\[
 \begin{gathered}
 z_m=P_mu^*,\qquad e_m=u^*-z_m,\qquad r_m=\cnorm{e_m}_*,\\
 S_m^{\rm fin}=\max_{d\in\cD_{\rm fin}}|a_*(e_m,d)|,\qquad
 \ell_m=\cnorm{(I-P_m)d_{m+1}}_*>0 .
 \end{gathered}
\]
Define $\theta_m=|a_*(e_m,d_{m+1})|/S_m^{\rm fin}$ if
$S_m^{\rm fin}>0$, and $\theta_m=0$ otherwise.
The vector $(I-P_m)d_{m+1}$ is the innovation of the candidate
$d_{m+1}$, that is, its component orthogonal to $\widehat\Phi_m$, and
$\ell_m$ is the norm of that innovation.
Let $Q_0=I-P_{m_0}$ and assume explicit coefficients satisfy
\begin{equation}
 v_B=\sum_{d\in\cD_{\rm fin}}a_dQ_0d,\quad
 \sum_d|a_d|\leq B_{\rm entry},\quad
 \cnorm{Q_0u^*-v_B}_*\leq\sigma,\qquad B_{\rm entry}\geq0.
 \label{eq:c5-entry-source}
\end{equation}
The vectors \(Q_0d\) retain their projected norms. The coefficient budget \(B_{\rm entry}\) belongs to this approximate restarted representation, while \(B_{\rm src}\) denotes the continuous atomic source budget.

\begin{theorem}[Weighted finite-trajectory source bound]
\label{thm:c5-source}
Under \eqref{eq:c5-spaces}--\eqref{eq:c5-entry-source}, take
$0\leq\underline\theta_m\leq\theta_m$ and set
\[
 W_j=\sum_{m=m_0}^{m_0+j-1}\frac{\underline\theta_m^2}{\ell_m^2},
 \qquad x_0=(r_{m_0}^2-\sigma^2)_+ .
\]
If $B_{\rm entry}>0$ and $x_0>0$, then for $0\leq j\leq J$,
\begin{equation}
 (r_{m_0+j}^2-\sigma^2)_+
 \leq\left(x_0^{-1}+\frac{W_j}{4B_{\rm entry}^2}\right)^{-1}.
 \label{eq:c5-source-initial}
\end{equation}
In particular, if $W_j>0$,
$r_{m_0+j}^2\leq\sigma^2+4B_{\rm entry}^2/W_j$.
If $x_0=0$ or $B_{\rm entry}=0$, all later errors are at most $\sigma$.
For $W_j=0$, projection monotonicity gives the bound.
\end{theorem}
\begin{proof}
For $w_m=(I-P_m)d_{m+1}$, Pythagoras gives
\begin{equation}
 r_m^2-r_{m+1}^2
 =\frac{|a_*(e_m,d_{m+1})|^2}{\ell_m^2}.
 \label{eq:c5-drop}
\end{equation}
The residual is orthogonal to both the starting and current spaces.
Consequently
\[
 r_m^2=a_*(e_m,Q_0u^*)\leq B_{\rm entry}S_m^{\rm fin}+\sigma r_m,
 \qquad
 S_m^{\rm fin}\geq\frac{(r_m^2-\sigma^2)_+}{2B_{\rm entry}}
\]
when $B_{\rm entry}>0$; the last step uses
$2(r_m^2-\sigma r_m)\geq r_m^2-\sigma^2$.
While both successive errors exceed $\sigma$, put
$x_m=r_m^2-\sigma^2>0$. Equation~\eqref{eq:c5-drop} implies
$x_{m+1}\leq x_m-b_mx_m^2$ with
$b_m=\underline\theta_m^2/(4B_{\rm entry}^2\ell_m^2)$.
Taking reciprocals proves \eqref{eq:c5-source-initial} up to a crossing.
After a crossing, monotonicity makes the positive part zero.
If $S_m^{\rm fin}=0$ or $B_{\rm entry}=0$, the source pairing gives
$r_m^2\leq\sigma r_m$ directly. This also covers zero errors.
\end{proof}

Barron--Cohen--Dahmen--DeVore prove in
\cite[Theorem~2.3, equation~(2.31)]{BarronCohenDahmenDeVore2008}
that exact OGA in a Hilbert space satisfies
\begin{equation}
 \cnorm{f-f_N}_H^2\leq\cnorm{f-h}_H^2+
            \frac{4\cnorm h_{L^1(\cD)}^2}{N},\qquad h\in L^1(\cD).
 \label{eq:c5-barron}
\end{equation}
Here \(L^1(\cD)\) denotes the atomic variation space. In the restarted setting, \(f=Q_0u^*\), \(h=v_B\), \(\cnorm{f-h}_*\le\sigma\), and the coefficient budget on the projected dictionary is at most \(B_{\rm entry}\). The source pairing and reciprocal estimate use this classical mechanism, with \Cref{thm:c5-source} retaining the observed weights \(\underline\theta_m^2/\ell_m^2\) along the generated spaces.
If $\underline\theta_m\geq\underline\theta>0$, then $\ell_m\leq1$
gives $W_j\geq j\underline\theta^2$.
The upper bound decreases as the accumulated weight $W_j$ grows.

Li--Siegel \cite[Corollary~4.2 and equation~(4.11)]{LiSiegelEntropy2024} obtain the approximate-source inequality
\[
 \frac12\big(\|r\|^2-\|f-h\|^2\big)
 \leq \|h\|_{\mathcal A}\,|\langle r,g\rangle|,
\]
which leads to the same innovation-weighted projection recursion. The
present theorem specializes this mechanism to the restarted, unrenormalized
projected dictionary and retains the observed
\(\underline\theta_m^2/\ell_m^2\) weights generated by the nonlinear selection rule.

\subsection{An exact quartic energy connection}
\label{sec:c5-quartic}

For the one-dimensional pure or regularized quartic energy define
\[
 X=\{u\in W^{1,4}(0,1):\textstyle\int_0^1u=0\},\qquad
 \cE_\varepsilon(u)=\int_0^1
       \left(\tfrac{\varepsilon^2}2|u'|^2+\tfrac14|u'|^4\right)\cdd x-\ell(u).
\]
An additive constant in the density is immaterial.
Assume $u^*\in X$ solves the weak equation and let $Y\subset X$
contain $u^*$ and the spaces/dictionary of
Section~\ref{sec:c5-source}. Use
\[
 a_*(v,w)=\int_0^1(\varepsilon^2+3|(u^*)'|^2)v'w'\cdd x .
\]
For $\varepsilon>0$ this is positive definite on $Y$.
For $\varepsilon=0$ assume explicitly that
$\int |(u^*)'|^2|v'|^2>0$ for every nonzero $v\in Y$;
the sufficient condition $(u^*)'\ne0$ almost everywhere ensures this
on the zero-mean space. Finiteness follows from H\"older's inequality.

For any $\widehat G_m\in\widehat\Phi_m$, define
\[
 D_m=\cnorm{\widehat G_m-z_m}_*,\quad
 v_m=\widehat G_m'-(u^*)',\quad
 R_m=\int_0^1\left((u^*)'v_m^3+\tfrac14v_m^4\right)\cdd x .
\]
\begin{proposition}[Quartic identity and comparison]
\label{prop:c5-quartic}
For every such state,
\begin{equation}
 \widehat\delta_m=\tfrac12(r_m^2+D_m^2)+R_m.
 \label{eq:c5-quartic}
\end{equation}
If $D_m\leq q_{\rm corr}r_m$ and
$|R_m|\leq\rho(r_m^2+D_m^2)$ throughout a finite range, where
$q_{\rm corr}\geq0$ and $0\leq\rho<1/2$, then
\begin{equation}
 (\tfrac12-\rho)r_m^2\leq\widehat\delta_m
 \leq C_Er_m^2,\qquad C_E=(\tfrac12+\rho)(1+q_{\rm corr}^2).
 \label{eq:c5-quartic-comparison}
\end{equation}
If \(U_m\) is a source upper bound for \(r_m^2\), the additive estimate is \(\widehat\delta_m\le U_m/2+D_m^2/2+|R_m|\).
\end{proposition}
\begin{proof}
Set $a=(u^*)'$ and $v=v_m$. Subtracting the linear part of the density
at $a$ gives
$\varepsilon^2v^2/2+3a^2v^2/2+av^3+v^4/4$.
The weak equation cancels the integrated linear part, including the load.
Also $\widehat G_m-z_m$ lies in the active space and
$u^*-z_m$ is orthogonal to it, so
$\cnorm{\widehat G_m-u^*}_*^2=D_m^2+r_m^2$.
This proves the identity; the bounds follow by substitution.
\end{proof}

In \eqref{eq:c5-quartic}, \(D_m\) and \(R_m\) account for the nonlinear correction at the computed state. The numerical records report the quadrature defect separately. The manufactured solution determines the values of \(q_{\rm corr}\), \(\rho\), and the source remainder used in the finite-trajectory comparisons.

\subsection{Projection bounds under a power-type decrease condition}
\label{sec:c5-power}

Retain the projections of Section~\ref{sec:c5-source} and set
\[
 x_m=r_m^2,\qquad
 \gamma_m=\frac{|a_*(e_m,d_{m+1})|}{\ell_m}.
\]
Then \eqref{eq:c5-drop} is $x_{m+1}=x_m-\gamma_m^2$.
When $S_m^{\rm fin}>0$, $\gamma_m=\theta_mS_m^{\rm fin}/\ell_m$;
otherwise $\gamma_m=0$ by its direct definition.
The power-type condition below relates the normalized innovation decrease to the projection residual.

\begin{theorem}[Explicit power-type finite-range bound]
\label{thm:c5-power}
Let $\alpha>0$, $c_0>0$, and assume, at every step with $r_m>0$
in $m_0\leq m<m_0+J$,
\begin{equation}
 \gamma_m^2\geq c_0r_m^{2+2/\alpha}.
 \label{eq:c5-power-condition}
\end{equation}
For $r_{m_0}>0$ and $0\leq j\leq J$,
\begin{equation}
 r_{m_0+j}^2\leq
 \left(r_{m_0}^{-2/\alpha}+\frac{c_0j}{\alpha}\right)^{-\alpha}.
 \label{eq:c5-power}
\end{equation}
If $r_{m_0}=0$, every later projection error is zero.
For the quartic models, if the relative conditions of
Proposition~\ref{prop:c5-quartic} also hold, then
\begin{equation}
 \widehat\delta_{m_0+j}\leq C_E
 \left(r_{m_0}^{-2/\alpha}+\frac{c_0j}{\alpha}\right)^{-\alpha}.
 \label{eq:c5-power-energy}
\end{equation}
Natural-distance and Sobolev bounds are obtained by taking respectively
the one-half and one-quarter powers, with the constants of
\CGAEnergyTheoremReference.
\end{theorem}
\begin{proof}
For $x_{m+1}>0$, convexity gives
\[
 x_{m+1}^{-1/\alpha}-x_m^{-1/\alpha}
 \geq\frac{x_m-x_{m+1}}{\alpha x_m^{1+1/\alpha}}
 \geq\frac{c_0}{\alpha}.
\]
Summation proves \eqref{eq:c5-power}. If a residual reaches zero,
nestedness gives the remaining cases. Substitute the projection bound
into \eqref{eq:c5-quartic-comparison}, then apply
\CGAEnergyTheoremReference.
\end{proof}

A sufficient factorization on positive-error steps is
\begin{equation}
 S_m^{\rm fin}\geq c_1r_m^{2\nu},\qquad
 \theta_m/\ell_m\geq c_2,\qquad
 \nu=\tfrac12+\tfrac1{2\alpha},\qquad c_1,c_2>0 .
 \label{eq:c5-power-factors}
\end{equation}
Indeed $\gamma_m^2\geq(c_1c_2)^2r_m^{4\nu}$ and
$4\nu=2+2/\alpha$. Thus $c_0=(c_1c_2)^2$ is admissible.
The residual visibility factor and the selected-direction factor are
the two factors in this lower bound.

For a finite sequence with strictly positive drops, every prescribed
$\alpha>0$ admits
\[
 c_{\min}(\alpha;J)=
 \min_{\substack{m_0\leq m<m_0+J\\x_m>0}}
 \frac{x_m-x_{m+1}}{x_m^{1+1/\alpha}}>0 .
\]
This constant defines an a posteriori envelope on the specified sequence. Continuation evaluates the same inequality with the calibration constants held fixed. The numerical records give the explicit bounds, constants, and index ranges.

\ifdefined\CGAChapterFiveEmbedded
\else
\begin{thebibliography}{9}
\bibitem{BarronCohenDahmenDeVore2008}
A.~R. Barron, A. Cohen, W. Dahmen, and R.~A. DeVore,
Approximation and learning by greedy algorithms,
\emph{Annals of Statistics} 36 (2008), 64--94.
\href{https://doi.org/10.1214/009053607000000631}{doi:10.1214/009053607000000631}.
\bibitem{LiSiegelEntropy2024}
Y. Li and J.~W. Siegel, Entropy-based convergence rates of greedy algorithms,
\emph{Mathematical Models and Methods in Applied Sciences} 34 (2024), 779--802.
\href{https://doi.org/10.1142/S0218202524500143}{doi:10.1142/S0218202524500143}.
The proof numbering used here is from
\href{https://arxiv.org/abs/2304.13332v2}{arXiv:2304.13332v2},
23 February 2024 (typeset 26 February).
\bibitem{DereventsovTemlyakov2022}
A. Dereventsov and V.~N. Temlyakov,
Biorthogonal greedy algorithms in convex optimization,
\emph{Applied and Computational Harmonic Analysis} 60 (2022), 489--511.
\href{https://doi.org/10.1016/j.acha.2022.05.001}{doi:10.1016/j.acha.2022.05.001}.
\bibitem{XuXu2025}
J. Xu and X. Xu,
Randomized greedy algorithms for neural network optimization in solving
partial differential equations,
\emph{Journal of Scientific Computing} 105 (2025), article 26.
\href{https://doi.org/10.1007/s10915-025-03050-5}{doi:10.1007/s10915-025-03050-5}.
\end{thebibliography}
\fi
\ifdefined\CGAChapterFiveStandalone
\def\CGAChapterFiveEndDocument{\end{document}}
\else
\let\CGAChapterFiveEndDocument\relax
\fi
\CGAChapterFiveEndDocument
